\documentclass[12pt]{article}
\usepackage[T1]{fontenc}
\usepackage[utf8]{inputenc}
\usepackage{lmodern}
\usepackage{textcomp}
\usepackage{enumitem}
\usepackage{geometry}
\geometry{margin=1in}
\begin{document}
\begin{center}
Answer Key - Sociology - Undergraduate Level Assessment 14 \
\small (Answers are ordered exactly as in the question paper.)
\end{center}

\section*{Multiple Choice}
\begin{enumerate}[label=\arabic*.]
\item \textbf{Answer:} D
\\ \textit{Options:} A) flexibility was reduced by the introduction of detailed daily worksheets; B) the state had greater control than managers over production processes; C) ownership was transferred from small shareholders to senior managers; D) employment depended on performance rather than trust and commitment
\\ \textit{Note:} Correct option: D - employment depended on performance rather than trust and commitment
\item \textbf{Answer:} C
\\ \textit{Options:} A) there are no clearly defined projects into which the money can be directed; B) the United Nations has refused to call on rich countries to provide it; C) debt repayments with interest can be greater than the amount of money received; D) debt repayments with interest can be greater than the amount of money received
\\ \textit{Note:} Correct option: C - debt repayments with interest can be greater than the amount of money received
\item \textbf{Answer:} D
\\ \textit{Options:} A) age; B) income; C) verbal fluency; D) occupation
\\ \textit{Note:} Correct option: D - occupation
\item \textbf{Answer:} A
\\ \textit{Options:} A) ways of acting, thinking, and feeling that are collective and social in origin; B) the way scientists construct knowledge in a social context; C) data collected about social phenomena that are proven to be correct; D) ideas and theories that have no basis in the external, physical world
\\ \textit{Note:} Correct option: A - ways of acting, thinking, and feeling that are collective and social in origin
\item \textbf{Answer:} C
\\ \textit{Options:} A) cults of the capital; B) capital culture; C) cultural capital; D) culpable capture
\\ \textit{Note:} Correct option: C - cultural capital
\end{enumerate}

\section*{True / False}
\begin{enumerate}[label=\arabic*.]
\setcounter{enumi}{5}
\item \textbf{Answer:} True
\\ \textit{Options:} A) True; B) False
\\ \textit{Note:} LLM-generated true statement based on MCQ.
\item \textbf{Answer:} False
\\ \textit{Options:} A) True; B) False
\\ \textit{Note:} LLM-generated false statement. Correct answer: 'community alternatives to imprisonment and institutional care'.
\item \textbf{Answer:} False
\\ \textit{Options:} A) True; B) False
\\ \textit{Note:} LLM-generated false statement. Correct answer: 'universalistic benefits and public sector employment'.
\item \textbf{Answer:} True
\\ \textit{Options:} A) True; B) False
\\ \textit{Note:} LLM-generated true statement based on MCQ.
\item \textbf{Answer:} True
\\ \textit{Options:} A) True; B) False
\\ \textit{Note:} LLM-generated true statement based on MCQ.
\end{enumerate}

\section*{Long Form Response}
\begin{enumerate}[label=\arabic*.]
\setcounter{enumi}{10}
\item \textbf{Answer:} Scapegoating. Explain why this is the correct answer and provide supporting evidence. Reference key facts or concepts that support this choice.
\\ \textit{Note:} Long-form derived from MCQ; award credit for stating the correct idea and giving supporting reasoning.
\item \textbf{Answer:} the exploitation process. Explain why this is the correct answer and provide supporting evidence. Reference key facts or concepts that support this choice.
\\ \textit{Note:} Long-form derived from MCQ; award credit for stating the correct idea and giving supporting reasoning.
\end{enumerate}

\vfill
\noindent\textit{End of Answer Key}
\end{document}